% FULL_SOLUTION_SCOPE: {"target":"Entire recorded catalogue statement: a closed partition expression for S_k of a sum of arbitrary n-by-n matrices over any field, 0<=N<=k<=n, with elementary symmetric functions of ordered words permitted.","result":"Theorem 3.1 and equation (6), with the self-contained proof of the primitive cyclic-word product in Lemma 2.2.","coverage":"All fields and characteristics, all matrix tuples including noncommuting and singular matrices, all 0<=N<=k<=n, empty sum N=0, empty partition k=0, N=1, and determinant case k=n. In fact the identity holds over every commutative ring.","dependencies":"Established Newton-Girard and Amitsur primitive cyclic-word identities, attributed to Amitsur (1980) and Reutenauer-Schutzenberger (1987); all identities needed for the answer are re-proved using elementary determinants, formal power series over a rational polynomial ring, cyclic-word orbit counting, and integral polynomial specialization.","human_verification":"None reported.","unverified":"The manuscript has no asserted mathematical gap, but all proofs remain without independent human or theorem-prover verification. Exact computations and same-model Prompt B review support only their stated scopes.","novelty":"Source and literature checked on 2026-10-12. This is a catalogue completion using established mathematics, with explicit partition organization and a self-contained proof; no claim of a new Amitsur or Newton-Girard identity. The complete 1987 primary paper was inspected; the 1980 full text was not accessible."}
\documentclass[11pt]{article}
\usepackage[margin=25mm]{geometry}
\usepackage[T1]{fontenc}
\usepackage{lmodern}
\usepackage{amsmath,amssymb,amsthm}
\usepackage{xurl,hyperref}
\hypersetup{colorlinks=true,urlcolor=blue,citecolor=blue,linkcolor=blue}
\setlength{\emergencystretch}{3em}
\newcommand{\tr}{\operatorname{tr}}
\newtheorem{theorem}{Theorem}[section]
\newtheorem{lemma}[theorem]{Lemma}
\title{A division-free partition formula for elementary symmetric functions\newline of a sum of matrices\thanks{Catalogue statement, ``Elementary symmetric of a sum of matrices.''}}
\author{OpenAI GPT 6.1 Sol (Max reasoning)\\[0.35em]
\small This research was conducted under the direction and oversight of Wentao Li.}
\date{12 October 2026}
\begin{document}
\section*{A division-free partition formula for elementary symmetric functions of a sum of matrices}
\noindent\textbf{OpenAI GPT 6.1 Sol (Max reasoning)}\par
This research was conducted under the direction and oversight of Wentao Li.\par
12 October 2026
\section*{Abstract}
We give a finite sum indexed by integer partitions of $k$ for
$S_k(A_1+\cdots+A_N)$, using only elementary symmetric functions of
ordered matrix words. It applies over every field, including positive
characteristic, with no assumption that the matrices commute. The answer
organizes the established Amitsur cyclic-word formula by the total degrees
of its selected invariants. A self-contained proof and exact checks are supplied.


\section{Original target and established context}
For $A\in M_n(F)$, where $F$ is any field, set
\[
 S_j(A)=\sum_{\substack{I\subseteq\{1,\ldots,n\}\\|I|=j}}
             \det A[I,I],\qquad S_0(A)=1.
\]
Thus $S_n(A)=\det A$, $S_1(A)=\tr A$, and
\begin{equation}\label{det-generating}
 \det(I_n-tA)=\sum_{j=0}^n(-1)^jS_j(A)t^j.
\end{equation}
For convenience set $S_j(A)=0$ for $j>n$. Determinants of zero-sized
matrices and empty products are $1$.

The recorded target~\cite{catalogue,costas} is to find a closed expression
organized by integer partitions for
\[
 S_k(A_1+\cdots+A_N),\qquad 0\leq N\leq k\leq n,
\]
in matrix quantities built from the $A_i$, permitting elementary symmetric
functions of their ordered words. Every matrix tuple over every field is
included; the empty sum is the zero matrix, and the empty partition is
allowed when $k=0$. No commutativity hypothesis is present.

Costas-Santos~\cite[\S\S1,5]{costas} gives low-degree formulas and asks for
a general closed expression involving partitions and elementary symmetric
functions of words. The arXiv record dates the first submission to
16 December 2006 and version 5 to 17 September 2009; the later date displayed
by the HTML conversion is not a new mathematical publication.

Amitsur's 1980 paper~\cite{amitsur} addresses characteristic coefficients
of sums in terms of coefficients of words. We consulted its publisher
abstract and bibliographic record, but could not access its full text.
The full primary paper of Reutenauer and Sch\"utzenberger~\cite{rs} was
visually inspected: its theorem on p.~300 and equations (6)--(7) on p.~301
give the cyclic-word formula over a commutative ring; its p.~302 addendum
credits Amitsur. The reference request~\cite{mo} was a bibliographic lead.
The present answer applies this established mathematics to the exact
recorded target and gives its partition organization; no novelty of the
underlying identities is claimed. This does not revise the catalogue's
original \texttt{open} source label.

\section{Primitive words and the integral product identity}
A word $w=(i_1,\ldots,i_d)$ in the alphabet $\{1,\ldots,N\}$ has length
$|w|=d$ and matrix value
$A_w=A_{i_1}\cdots A_{i_d}$, in precisely that order.
A nonempty word is \emph{primitive} if it is not a power $v^r$ with $r>1$.
Identify words only under cyclic rotation, and choose the lexicographically
least representative of each primitive cyclic class. Let $\mathcal C_N$
be these representatives and
$\mathcal C_{N,k}=\{w\in\mathcal C_N:|w|\leq k\}$.
These sets depend only on the alphabet, not on the matrix entries.
When $N=0$ there are no nonempty words.

\begin{lemma}[Cyclic invariance]\label{cyclic}
For matrices $X,Y$ of the same size over any commutative ring $R$,
$\det(I-uXY)=\det(I-uYX)$ in $R[u]$. Consequently $S_j(A_w)$
is independent of the cyclic representative of $w$.
\end{lemma}
\begin{proof}
Eliminate either off-diagonal block in
$\left(\begin{smallmatrix}I&uX\\Y&I\end{smallmatrix}\right)$
by multiplication with a block triangular matrix with identity diagonal
blocks. The two resulting determinants are $\det(I-uXY)$ and
$\det(I-uYX)$. Compare coefficients using~\eqref{det-generating}.
Every one-step cyclic rotation has the form $XY\mapsto YX$.
\end{proof}

\begin{lemma}[Amitsur product, with integral specialization]\label{product}
For any finite tuple of $n\times n$ matrices over a commutative ring $R$,
\begin{equation}\label{amitsur-product}
 \det\left(I_n-t\sum_{i=1}^N A_i\right)
       =\prod_{w\in\mathcal C_N}\det(I_n-t^{|w|}A_w)
       \quad\hbox{in }R[[t]].
\end{equation}
The product is formal: its coefficients through degree $k$ depend only
on the finite set $\mathcal C_{N,k}$.
\end{lemma}
\begin{proof}
First use generic matrices $X_1,\ldots,X_N$ whose entries are distinct
indeterminates, and let $P=\mathbb Z[x_{i,a,b}]$ and
$Q=\mathbb Q[x_{i,a,b}]$. Put $B=\sum_iX_i$.
For any matrix $M$ over $Q$, the scalar formal series
$f(t)=\det(I-tM)$ has constant term $1$. Differentiating the determinant
by multilinearity, and using the adjugate identity, gives
\[
 \frac{f'(t)}{f(t)}=-\tr\bigl(M(I-tM)^{-1}\bigr)
                 =-\sum_{m\geq1}\tr(M^m)t^{m-1}.
\]
Termwise integration with zero constant term yields
\begin{equation}\label{log-determinant}
 \log\det(I-tM)=-\sum_{m\geq1}\frac{\tr(M^m)}{m}t^m.
\end{equation}
All operations here are formal over a rational algebra.

Expand $B^m$ as the sum of all ordered words of length $m$.
The cyclic orbit of such a word has size $d$ for a unique divisor $d$ of
$m$: the rotations fixing the word form a subgroup of
$\mathbb Z/m\mathbb Z$, so the least positive period $d$ divides $m$.
The word is the $(m/d)$-th power of a primitive word of length $d$.
Each primitive cyclic class of length $d$ produces exactly one such orbit,
with exactly $d$ distinct linear words. Their traces agree by cyclic
invariance of trace. Therefore
\begin{equation}\label{orbit-trace}
 \tr(B^m)=\sum_{d\mid m}d
                \sum_{w\in\mathcal C_N,\ |w|=d}
                       \tr(X_w^{m/d}).
\end{equation}
Using $m=dr$ in~\eqref{log-determinant} gives
\[
 \begin{aligned}
 \log\det(I-tB)
   &=-\sum_{w\in\mathcal C_N}\sum_{r\geq1}
                    \frac{\tr(X_w^r)}r t^{r|w|}\\
   &=\sum_{w\in\mathcal C_N}\log\det(I-t^{|w|}X_w).
 \end{aligned}
\]
Each coefficient receives only finitely many contributions. Formal
exponentiation proves~\eqref{amitsur-product} over $Q$.

Both sides of that identity have coefficients in $P$: each right-hand
factor is a determinant with integer polynomial coefficients, and only
finitely many factors affect a fixed degree. Since $P\hookrightarrow Q$
is injective, equality already holds coefficient by coefficient in $P$.
Now specialize the indeterminates to the entries of any given matrices
over $R$ via the homomorphism $P\to R$. This proves the identity over $R$,
including finite fields and rings of positive characteristic. If $N=0$,
both sides are $1$, with an empty product on the right.
\end{proof}

\section{The finite partition expression}
Write $\lambda\vdash k$ for an integer partition
$\lambda=(\lambda_1\geq\cdots\geq\lambda_\ell>0)$ of $k$.
For each such partition let $D_\lambda(N,n)$ be the following finite,
explicitly enumerable set of tuples:
\begin{equation}\label{decorations}
 \begin{split}
 D_\lambda(N,n)=\bigl\{((w_1,r_1),\ldots,(w_\ell,r_\ell)): &
 w_h\in\mathcal C_{N,k},\quad 1\leq r_h\leq n,\quad
 r_h|w_h|=\lambda_h;\\
 &w_h\ne w_j\ (h\ne j);\\
 &w_h<w_{h+1}\text{ whenever }\lambda_h=\lambda_{h+1}\bigr\}.
 \end{split}
\end{equation}
The last comparison uses lexicographic order and fixes the order inside
each block of equal parts. For the empty partition, $D_\varnothing(N,n)$
contains the empty tuple. An empty inner sum is $0$.

\begin{theorem}[Closed partition formula]\label{main}
For every field $F$, every $A_1,\ldots,A_N\in M_n(F)$, and
$0\leq N\leq k\leq n$, one has
\begin{equation}\label{partition-formula}
 \boxed{\displaystyle
 S_k\left(\sum_{i=1}^N A_i\right)=
 \sum_{\lambda\vdash k}
 \sum_{((w_h,r_h))_{h=1}^{\ell}\in D_\lambda(N,n)}
       (-1)^{\,k+\sum_{h=1}^{\ell}r_h}
       \prod_{h=1}^{\ell}S_{r_h}(A_{w_h}).}
\end{equation}
The same formula holds for every finite $N$ and $k\geq0$ over any
commutative ring, with the convention $S_k=0$ for $k>n$.
\end{theorem}
\begin{proof}
In the product~\eqref{amitsur-product}, the factor for $w$ is
\[
 1+\sum_{r=1}^n(-1)^r S_r(A_w)t^{r|w|}.
\]
To obtain degree $k$, choose a nonconstant term from finitely many
\emph{distinct} factors, each supplying an index $r_h$. The selected
degrees $r_h|w_h|$ have sum $k$. Sort them into an integer partition
$\lambda$, and sort the representatives within equal-degree blocks.
This is a bijection between product selections and the tuples
in~\eqref{decorations}; no factorial or extra cyclic-orbit multiplicity
is needed. Conversely every such tuple specifies one selection.
Its coefficient is $(-1)^{\sum_h r_h}\prod_h S_{r_h}(A_{w_h})$.
The degree-$k$ coefficient on the left is $(-1)^k S_k(\sum_iA_i)$,
so multiplication by $(-1)^k$ proves~\eqref{partition-formula}.
\end{proof}

This is a closed finite expression rather than a recurrence: for each
part $s$, enumerate primitive representatives of lengths $d\mid s$ with
$d\leq k$, then take $r=s/d\leq n$ and enforce the displayed distinctness
and ordering rules. Every scalar invariant in the answer is an $S_r$
of an ordered word in the input matrices. Each part records its total
degree $r|w|$ in those matrices. Cyclic rotation is allowed;
reversal and other permutations of letters are not identified.

\paragraph{Boundary cases.}
For $k=0$ the only partition and decoration are empty, so the answer is
$1$; within the recorded range this forces $N=0$.
For $N=0$ and $k>0$ every decoration set is empty, giving $0$.
For $N=1$ the only primitive representative is the single letter $1$:
only $\lambda=(k)$ contributes when $k>0$, giving $S_k(A_1)$ with sign
$(-1)^{2k}=1$. The case $k=n$ is the required determinant formula.
If $n=0$ is permitted, the sole recorded case is again $N=k=0$.
Singular matrices and zero input matrices need no separate hypothesis.

\section{The characteristic-zero trace partition formula}
For comparison, when $F$ has characteristic zero define
\[
 p_j=\sum_{(i_1,\ldots,i_j)\in\{1,\ldots,N\}^j}
                S_1(A_{i_1}\cdots A_{i_j})
     =\tr\left(\left(\sum_i A_i\right)^j\right),
 \qquad
 z_\lambda=\prod_{j\geq1}j^{m_j(\lambda)}m_j(\lambda)!.
\]
Here $m_j(\lambda)$ is the multiplicity of the part $j$ and
$\ell(\lambda)=\sum_jm_j(\lambda)$. The familiar Newton--Girard
partition expression is
\begin{equation}\label{trace-partition}
 S_k\left(\sum_i A_i\right)=
   \sum_{\lambda\vdash k}
    \frac{(-1)^{k-\ell(\lambda)}}{z_\lambda}
                  \prod_{j\geq1}p_j^{m_j(\lambda)}.
\end{equation}
Indeed, substitute $-t$ in~\eqref{log-determinant} and expand
\[
 \det(I+tB)=\prod_{j\geq1}
          \exp\left(\frac{(-1)^{j-1}p_j}{j}t^j\right).
\]
Choosing multiplicity $m_j$ supplies the denominator $j^{m_j}m_j!$ and
sign $(-1)^{\sum_j(j-1)m_j}=(-1)^{k-\ell(\lambda)}$.
This proves~\eqref{trace-partition} without commuting any matrix letters.
Equation~\eqref{partition-formula} supplies the arbitrary-characteristic
answer. For example, over $\mathbb F_2$ the matrices $I_2$ and $0_2$
have identical traces of every positive power, but their $S_2$ values are
$1$ and $0$. Ordinary traces of words alone cannot distinguish this
one-letter example; higher $S_r$ are essential.

\section{Examples and checks of order, signs, and repetition}
For $n\geq2$ the degree-two specialization is
\begin{equation}\label{degree-two}
 S_2\left(\sum_iA_i\right)
   =\sum_i S_2(A_i)+\sum_{i<j}
          \bigl(S_1(A_i)S_1(A_j)-S_1(A_iA_j)\bigr).
\end{equation}
For $A=E_{12}$ and $B=E_{21}$ in dimension $2$, $AB=E_{11}$ and
$BA=E_{22}$, so the pair does not commute. Equation~\eqref{degree-two}
gives $S_2(A+B)=-1$, agreeing with the determinant of
$\left(\begin{smallmatrix}0&1\\1&0\end{smallmatrix}\right)$;
over $\mathbb F_2$ this value is $1$.

For two letters the primitive representatives through length $4$ are
\[
 A,\ B;\quad AB;\quad AAB,\ ABB;\quad AAAB,\ AABB,\ ABBB.
\]
In particular $ABAB$ is periodic, and is accounted for by $S_2(AB)$
in the factor belonging to $AB$. For $n\geq4$ the formula gives
\begin{equation}\label{degree-four}
\begin{aligned}
 S_4(A+B)={}&\sum_{\substack{a,b\geq0\\a+b=4}}S_a(A)S_b(B)
       -S_1(AB)\sum_{\substack{a,b\geq0\\a+b=2}}S_a(A)S_b(B)\\
 &+\bigl(S_1(A^2B)+S_1(AB^2)\bigr)\bigl(S_1(A)+S_1(B)\bigr)\\
 &-S_1(A^3B)-S_1(A^2B^2)-S_1(AB^3)+S_2(AB).
\end{aligned}
\end{equation}
This recovers the source's fourth-degree invariant language in every
characteristic.

Order beyond cyclic rotation matters already for $N=k=n=3$.
Take $A_1=E_{12}$, $A_2=E_{23}$, $A_3=E_{31}$.
Then $S_1(A_1A_2A_3)=1$ and $S_1(A_1A_3A_2)=0$.
These are distinct primitive cyclic classes. In degree $3$ they are
the only possibly nonzero terms of~\eqref{partition-formula}; their sum
is $1$, which equals $S_3(A_1+A_2+A_3)$ over every field.

\section{Verification, scope, and contributions}
\paragraph{Exact computation.}
Download the \href{https://api.github.com/repos/mehmetmars7/Erdosproblems-llm-hunter/git/blobs/ac0c2c73c5709a5b519e85791c2d0b804d1f83e6}{verified source blob}, decode its base64 \texttt{content} field, and save the unchanged bytes as \texttt{verifier.py}. Run \texttt{python3 verifier.py}. The checker uses Python 3.13.2, its standard library, and seed $3027$.
Its oracle sums principal minors using the Leibniz determinant, separately
from both the decorated-partition sum and the primitive-word product.
All $2293$ coefficient cases and $469$ auxiliary checks passed.
The comparisons cover integer examples through dimension $6$ and every
$0\leq N\leq n$; rational examples; finite fields
$\mathbb F_2,\mathbb F_3,\mathbb F_5,\mathbb F_7$ through dimension $5$;
every one of the $256$ ordered pairs of $2\times2$ matrices over
$\mathbb F_2$; and small generic polynomial pairs in dimensions $2$ and
$3$ with respectively $8$ and $18$ independent variables.
Additional examples over $\mathbb Z/4,\mathbb Z/6,\mathbb Z/8$ test ring
specialization. Some checks extend past degree $n$, where coefficients
must vanish. The auxiliary checks cover primitive-word counts, periodic
trace multiplicities, cyclic invariance, and the trace obstruction.
The program also checks~\eqref{trace-partition} over the rationals.

\paragraph{Full-solution scope declaration.}
Theorem~\ref{main} answers the entire recorded target: an explicit finite
integer-partition expression in the permitted invariants, for every matrix
tuple, every field, and every $0\leq N\leq k\leq n$, including all boundary
cases. Lemmas~\ref{cyclic} and~\ref{product} prove the needed identities;
the proof uses only elementary determinants, formal series, finite cyclic
orbit counting, and specialization of integer polynomials. The established
Amitsur and Newton--Girard machinery is attributed and re-proved here.
There is no asserted remaining mathematical gap for this target.
The separate positivity conjecture in the source is not part of catalogue
the recorded question and is not addressed by this answer.

\paragraph{Process and attribution.}
OpenAI GPT 6.1 Sol (Max reasoning) generated the proof, manuscript, and
verification code. Research directed and overseen by Wentao Li.
No other model was used. Filled mandatory Prompts A and B, draft and audit
snapshots, source-access notes, and the exact output report are retained in
the local research directory,
outside this repository. The checker is included with this contribution;
those research notes are not asserted to have been publicly posted.
The audit is a same-model review. No independent human proof verification
or theorem-prover formalization has been reported; the proofs remain
unverified in those senses. Computation does not replace the general proof.

\paragraph{Catalogue provenance.}
The target and identity metadata derive from ULAM.AI and the UnsolvedMath
Contributors~\cite{catalogue}, under CC BY 4.0. This separately attributed
attempt adds the formula, proof, and checks; it preserves the imported
target and source status. Source/literature and repository collision checks
were performed on 12 October 2026. No earlier attempt or matching pull
request was found in the checked snapshot and searches; this is not a claim
of an exhaustive historical literature search or new mathematical priority.

\section{Completion Estimate}
\noindent\textbf{COMPLETION: 100\%}
Theorem~\ref{main} supplies the requested closed partition expression for
all recorded fields, matrix tuples, parameter values, and boundary cases.
This is a full-solution claim for the recorded target using established
mathematics, subject to the verification limits stated above.

\begin{thebibliography}{9}
\bibitem{catalogue}
ULAM.AI and the UnsolvedMath Contributors,
\emph{UnsolvedMath: A Curated Collection of Open Mathematics Problems},
catalogue entry.
\url{https://huggingface.co/datasets/ulamai/UnsolvedMath}.
CC BY 4.0: \url{https://creativecommons.org/licenses/by/4.0/}.
Recorded statement: \textit{catalogue statement},
repository snapshot \texttt{1397a8c7f75e48a6343b0afe338390b365d344df}.

\bibitem{costas}
R. S. Costas-Santos,
\emph{On the elementary symmetric functions of a sum of matrices},
arXiv:math/0612464v5 (17 September 2009), especially \S\S1,5.
\url{https://arxiv.org/abs/math/0612464};
\url{https://arxiv.org/html/math/0612464v5}.

\bibitem{amitsur}
S. A. Amitsur,
\emph{On the characteristic polynomial of a sum of matrices},
Linear and Multilinear Algebra \textbf{8} (1980), 177--182.
\url{https://doi.org/10.1080/03081088008817315}.
Publisher abstract consulted; full text not accessed.

\bibitem{rs}
C. Reutenauer and M.-P. Sch\"utzenberger,
\emph{A formula for the determinant of a sum of matrices},
Letters in Mathematical Physics \textbf{13} (1987), 299--302.
\url{https://doi.org/10.1007/BF00401158}.
Primary scanned copy, all four pages inspected:
\url{https://igm.univ-mlv.fr/~berstel/Schutzenberger/Travaux/A/1987-2FormulaDeterminantMathPhys.pdf}.

\bibitem{mo}
P. Hartwig, \emph{Reference for a formula expressing the characteristic
polynomial of a sum of endomorphisms}, MathOverflow question 70231
(13 July 2011), with E. Formanek's answer (15 July 2011).
\url{https://mathoverflow.net/questions/70231}.
Consulted as a reference lead, not used as proof of the identity.
\end{thebibliography}
\end{document}
